
\maketitle


\section{Introduction}

Although the game ``rock, paper, scissors'' is often used as a random decider,
similar to a coin, it is in fact possible to win more consistently than
randomness would predict. This project uses data mining to construct an agent
that learns from prior games with its opponent, and then uses that information
to win games through predicion.


Our methodology involves steady learning, and near-continous
retraining in order to achieve the effect of an opponent that learns
as it plays. Part of the motivation and significance of this project
comes as an attempt to understand the processes people use to make decisions
where opponent prediction and ``randomness'' are rewarded factors.


\section{Data Collection}

There is no pre-constructed dataset that we can use for this project. Instead,
we will be collecting our own data by recording series' of rock paper scissors
games, both between two humans, and between a single human and the model. For
each game, we will collect the game id, which item each opponent played (one of
rock, paper, or scissors), and the result.



\section{Technical Approach}
The predicion model will take an input of every previously played game in the
session, and use the combination of the user input and our model's previously
generated predicted input, to determine it's next response. For this, we plan,
so far, to use a combination of a simple ``most common pick'' (uses the most
picked item), and an implementation of the the apriori algorithm to attempt to
find patterns in the user's previous responses.


There is a slight challenge of always collecting data, in that you need to pick
a time to train the model. The first idea for this was to recalculate our Apriori
and ``most common pick'' every few runs of the algorithm, although with how
light our data is, it is also possible to do it every run.

\section{Evaluation Methodology}
It should be fairly simple to evaluate the success of the model. The main
variable we care about is simply the win rate --- how often the model wins its
game. We will also look at the number of games it takes to start winning more
consistently; we want it to take fewer games to ``warm up''.


% win rate
% total wins
% games it takes to get to a good win rate


\section{Results and Discussion}
We don't have any results yet, since we have not begun work on the project.
